\MajorSection{النص المعتمد}

\textenglish{Richard Francis Burton. The Kasîdah of Hâjî Abdû El Yezdî}. نص العمل المنشور سنة ١٨٨٠، في النسخة الرقمية من مشروع غوتنبرغ، الكتاب رقم ٦٠٣٦. تتضمن النسخة المقدمة، والأقسام التسعة، والملاحظتين، والخاتمة. تاريخ الاطلاع ٣٠ أيلول ٢٠٢٦.

\href{https://www.gutenberg.org/ebooks/6036}{صفحة الكتاب في مشروع غوتنبرغ}

\href{https://www.gutenberg.org/cache/epub/6036/pg6036-images.html}{النص الإنجليزي الكامل المعتمد}

للتعريف بالإطار الأدبي للكتاب: \textenglish{The Kasidah}، موسوعة ويكيبيديا، مع إحالاتها إلى العمل وطبعاته.

\href{https://en.wikipedia.org/wiki/The\_Kasidah}{تعريف بالإطار الأدبي للقصيدة}

ولسياق مصطلح المادة المتألقة: \textenglish{William Crookes. On Radiant Matter}. محاضرة ألقيت في شيفيلد في ٢٢ آب ١٨٧٩، ونُشرت في \textenglish{Popular Science Monthly}، المجلد ١٦، تشرين الثاني ١٨٧٩. النص متاح في ويكي مصدر.

\href{https://en.wikisource.org/wiki/On\_Radiant\_Matter}{محاضرة وليام كروكس عن المادة المتألقة}

\MinorSection{مصدر التقديم عن الترجمة الموهومة}

\begin{english}
Glyn Pursglove (2011). Fakery, Serious Fun and Cultural Change: Some Motives of the Pseudo-Translator. Hermēneus: Revista de Traducción e Interpretación, 13.
\end{english}

وللمناقشة الخاصة ببرتون، ص ٤–٥.

\href{https://www5.uva.es/hermeneus/hermeneus/13/arti06_13.pdf}{دراسة غلين بورزغلوف عن دوافع المترجم الموهوم}
